\section{Related Work} \label{sec:soa}
The study of human activity has been an important goal of computer vision since its inception and has developed considerably in recent years. To address this problem, researchers have proposed several methods over the past years based on traditional Artificial Neural Networks and, more recently, on deep learning architectures. 


%###Se ha aplicado a diferentes campos...

%%In the current works, researchers have tried to summarize in groups of knowledge on the subject, among the main topics used to analyze human behavior are Human Behavior Analysis (HBA)\cite{chaaraoui2014vision}, Activities of Daily Living (ADL)\cite{wang2019deep} type applications that generally use data obtained with different types of sensors. Acceleration and angular velocity are modified according to the movements of the human body, so some human activities can be inferred. These sensors can usually be found in smart phones, watches, bands, glasses and helmets, Human Activity Recognition (HAR)\cite{jalal2019multi}, is common in all aspects of daily life. Therefore, this topic has become interesting for researchers, basic daily activities such as sitting, walking, running, and can be monitored with a high level of accuracy if the user carries a large number of sensor nodes mainly in smart phones. In our study the source of data to detect human activities is a video, Ambient Assisted Living (AAL)-type applications aim to improve the quality of life and maintain independence especially for the elderly and vulnerable using technology. The latest generation of AAL technologies uses many sensors integrated into the environment or in the home of people. These sensors are usually invasive and often activated by false alarms, to this an alternative option is the use of cameras that can detect mainly falls of people, as this type of abnormal activity is frequent in older people, for this reason there is a growing trend in solutions based on video and artificial vision for applications AAL\cite{chaaraoui2012}.

%# Soluciones parciales...
The works use different strategies to address partial problems of the challenge of analysing group behaviour. In this way, the works propose partial solutions to determine movements, actions, activities of groups, other focus on small groups or crowds, etc. In \cite{devanne2019recognition} the authors propose a vision-based solution to identify Activities of Daily Living (ADL), through skeletal data captured with an RGB-D camera. After the decomposition of a skeletal sequence into short time segments, the activities are classified through a two-layer network called Long-Short Term Memory Network (LSTM), which allows analyzing the sequence at different levels of temporal granularity. The proposal is evaluated in the dataset Watch-n-Patch\cite{DBLP:journals/corr/WuZSSSS16}, in which there are examples 11 different daily activities of people such as: bringing things from the fridge, turning things to the fridge, spilling liquid, drinking a liquid, leaving the kitchen, bringing things from the oven, putting things into the microwave, preparing food, filling a kettle, connecting a kettle to the electrical outlet, moving the kettle. The main contribution of the authors is a model of activities of multiple scales and time dependence, based on the comparison of the characteristics of the context that characterize the results of previous recognitions, and a hierarchical representation with a recognition layer of low level behavior units and another high level unit. That is, it is a solution that handles two different levels of semantics.

Human actions in video sequences are usually defined as three-dimensional (3D) space-time signals that characterize both the visual and dynamic appearance of the movement of the human beings and objects involved. Taking into account the success generated by the positive results of Convolutional Neural Networks (CNN) for image classification, recent attempts have been made to learn CNN 3D to recognize human actions in videos, but due to the high complexity of the training of 3D convolution cores and the need for large amounts of training videos that this type of networks requires for their training, there are only few success stories, being a broad topic still requiring maturity in research. 

Regarding the study of groups of people, advances in analysing behaviour are limited to very concrete and simple activities or actions, usually of short duration (low semantic component) such as a actions in sport games \cite{rahmad2018survey,chen2021lstm,muhammad2021human,ullah2021attention}, detection interactions of people inside a group \cite{wu2021comprehensive, deng2016structure,wang2017recurrent}, inter-group violence \cite{su2021deep, yao2021survey,ye2021campus}, among others. If we increase the number of people in the group, becoming crowds, the level of semantics is even lower, being specifically limited to tasks such as counting people and calculating crowd density \cite{ying2021survey, jingying2021survey,fan2022survey,bendali2021recent} or detecting movements of a mass of people or crowd collisions \cite{hao2019effective, nayan2019detecting, shehab2019statistical, zhang2019detection}, mainly for the purpose of security tasks. It is important to highlight the work in \cite{VAHORA201947}, in which the authors present a model of learning based on contextual relationships that uses a deep neural network to recognize activities in a video sequence. The proposed model involves contextual learning using a bottom-up approach, learning from individual human actions to group-level activities, and learning from scene information. Taking into account that it can identify group and individual activities would be considered a progress in this field of research, since it is one of the first works that can identify behaviors according to the number of people in the scene, clarifying that it is not yet reached the level of identification of behaviors in crowds with the same application \cite{VAHORA201947,vishwakarma2015human}.

%MC y OC
Finally, it is important to analyse the one-class classification approaches. They become particularly relevant when an adequate dataset is not available for all the classes that the classification models have to learn. There exists many situations in which it is necessary to identify one class among others but only having examples of this class. For example, in video surveillance, it is common to have information on normal actions or activities while it is very difficult to have examples of criminal actions. The same is valid for credit card transactions, where the large volume is related to legal transactions but the objective is to detect fraudulent ones. The models have to be trained without or scarce samples of abnormal situations, in which the goal is to learn from data the meaning of “normal”. Deviations or data different from this definition are considered as anomalies or “abnormal”. The problem of having most (or all) examples of a particular class becomes a bigger problem when using deep learning techniques as they are large consumers of data. Some proposals have been presented to address the one-class classification problem using deep neural networks. In particular, it is important to highlight the work of Chalapathy et al. \cite{chalapathy2019deep} that proposed a model of a one-class neural network (OC-NN) to detect anomalies in complex datasets. Among the work specifically addressing activity recognition in groups or crowds, different neural network models have been used to solve the problem. For example, the use of Convolutional Neural Networks (CNN) is shown in the work of Li et al. \cite{li2019object} where a new colorization of images including other information as optical flow is used as input of CNNs to detect objects and their anomalies. Also, in \cite{su2021deep}, Su et al. integrate the one-class Support Vector Machine into a CNN proposing the Deep One-Class (DOC) model. One widely used model has been autoencoders (AE) where they attempt to extract features from images to form a new space in which to decide the existence of normal activities. In this way, Saokrou et al. \cite{sabokrou2017fast} propose a cubic-patch-based method based on a cascade of AE to represent the information in the patches; Vu et al. \cite{vu2019robust} propose representation learning using Denoising Autoencoders (DAEs); and the works in \cite{xu2015learning,xu2017detecting} are based on multiple Stacked Denoising AutoEnocders (SDAEs). Finally, Generative Adversial Networks (GAN) trained using normal data have been used to learn an internal representation of the scene normality \cite{ravanbakhsh2017abnormal,ravanbakhsh2019training}.
